\documentclass[a4paper,11pt]{article}
\usepackage[utf8]{inputenc}
\usepackage[T1]{fontenc}
\usepackage[french]{babel}
\usepackage{amsmath}
\usepackage{amssymb}
\usepackage[francais,bloc,completemulti]{automultiplechoice}

\begin{document}

% \title{Automultiplechoice Example}
% \author{Your Name}
% \date{\today}
% \maketitle

\exemplaire{1}{
%%% debut de l’en-tête des copies :
\noindent{\bf QCM \hfill TEST}
\vspace*{.5cm}
\begin{minipage}{.4\linewidth}
\centering\large\bf Test\\ Examen du 01/01/2008\end{minipage}
\champnom{\fbox{
\begin{minipage}{.5\linewidth}
Nom et prénom :
\vspace*{.5cm}\namefielddots
\vspace*{1mm}
\end{minipage}
}}
\begin{center}\em
Durée : minutes.

Aucun document n’est autorisé.
L’usage de la calculatrice est interdit.
Les questions faisant apparaître le symbole \multiSymbole{} peuvent
présenter zéro, une ou plusieurs bonnes réponses. Les autres ont
une unique bonne réponse.
Des points négatifs pourront être affectés à de \emph{très
mauvaises} réponses.
\end{center}
\vspace{1ex}
%%% fin de l’en-tête

\begin{question}{Q1}
What is the capital of France?
\begin{reponses}
    \bonne{René Coty}
    \mauvaise{Alain Prost}
    \mauvaise{Marcel Proust}
    \mauvaise{Claude Monet}
\end{reponses}
\end{question}

\begin{question}{Q2}
What is the square root of 16?
\begin{reponses}
    \bonne{René Coty}
    \mauvaise{Alain Prost}
    \mauvaise{Marcel Proust}
    \mauvaise{Claude Monet}
\end{reponses}
\end{question}

\begin{question}{Q3}
Solve the equation $2x + 3 = 9$.
\begin{reponses}
    \bonne{René Coty}
    \mauvaise{Alain Prost}
    \mauvaise{Marcel Proust}
    \mauvaise{Claude Monet}
\end{reponses}
\end{question}

\begin{questionmult}{pref}
   Parmi les villes suivantes, lesquelles sont des préfectures~?
   \begin{reponses}
   \bonne{Poitiers}
   \mauvaise{Sainte-Menehould}
   \bonne{Avignon}
   \end{reponses}
   \end{questionmult}

\AMCaddpagesto{1}
}
\end{document}